% ECONOMICS_PROBLEM: {"id": "I13-393", "title": "Heterogeneous moral hazard and optimal insurance", "jel_code": "I13", "source_id": "OP-0052"}
% !TeX program = xelatex
\documentclass[11pt,a4paper]{article}
\usepackage[margin=23mm]{geometry}
\usepackage{fontspec}
\setmainfont{Latin Modern Roman}
\usepackage{amsmath,amssymb,mathtools}
\usepackage{xcolor,enumitem,booktabs,longtable,array,xurl,hyperref}
\definecolor{muted}{HTML}{53636C}
\hypersetup{colorlinks=true,urlcolor=blue,linkcolor=blue}
\setlength{\parindent}{0pt}
\setlength{\parskip}{5pt}
\newcommand{\R}{\mathbb R}
\newcommand{\E}{\mathbb E}
\newcommand{\N}{\mathbb N}
\newcommand{\ind}{\mathbf 1}
\newcommand{\Var}{\operatorname{Var}}
\newcommand{\Cov}{\operatorname{Cov}}
\newcommand{\supp}{\operatorname{supp}}
\DeclareMathOperator*{\argmin}{arg\,min}
\DeclareMathOperator*{\argmax}{arg\,max}
\newcommand{\entrymeta}[3]{{\sffamily\small\color{muted}\textbf{\texttt{#1}}\quad|\quad #2\par #3\par}}
\newcommand{\Problemref}[1]{\texttt{#1}}
\newcommand{\JELref}[2]{\texttt{#2} (JEL \texttt{#1})}
\begin{document}
\section*{Heterogeneous moral hazard and optimal insurance}
\textbf{Problem ID:} \texttt{I13-393}\quad \textbf{JEL:} \texttt{I13}\par
% BEGIN ECONOMICS STATEMENT
\label{problem:OP-0052}
\label{atlas:prob:H2}
\noindent\textbf{Original atlas identifier:} H2; entry 52. \textbf{Field:} Health economics. \textbf{Original tractability label:} Medium.

\noindent\textbf{Classification:} Parameterized theoretical/empirical research agenda; not a certified unsolved theorem.

\subsection*{Definitions and assumptions}
Let population type $\theta$ include health risks, wealth, treatment benefit, preferences, and utilization responsiveness. An insurance contract $c$ specifies a premium, deductible, treatment-specific coinsurance, out-of-pocket cap, and covered services; patient information legally usable for pricing is fixed. A model $m$ specifies contract choice, nonlinear budget constraints, treatment decisions, provider behavior, and health outcomes. The population includes patients and all affected financing taxpayers. Let $U_\theta(c,\mathcal C;m)$ be expected cardinal utility, including consumption risk, health, and the incidence of the stated financing taxes, and $w(\theta)\geq0$ social weights. Let $G(\mathcal C;m)$ be subsidies plus administration minus the financing tax receipts already modeled in private budgets. The parameter $\lambda$ values only residual public-budget opportunity costs in utility units; it is not a second tax-incidence charge. A feasible menu $\mathcal C$ must satisfy insurer budget balance/participation, legal nondiscrimination rules, and incentive-compatible self-selection where private types choose contracts.

\subsection*{Formal research question}
\emph{Editorial formulation: the mathematical target below is not a theorem quoted from the references.}

Characterize menus maximizing
\[
 \sup_{\mathcal C\in\mathfrak C}\left\{\int w(\theta)U_\theta(c^*(\theta,\mathcal C),\mathcal C;m)\,dF(\theta)-\lambda G(\mathcal C;m)\right\},
\]
Here $\mathfrak C$ is the feasible menu family, $F$ the population-type distribution, and $c^*$ the selected equilibrium contract (including a null contract for nonbuyers). Utility depends on the whole menu through its financing taxes. Determine what exogenous contract or price variation identifies the joint distribution of health risk and treatment-specific price responses, including correlation between them. Recover or bound the health loss from discouraged care and the risk-protection gain from lower cost sharing. Report $\varepsilon$-optimal feasible menus for models observationally equivalent in aggregate expenditure data; a spending response alone does not determine whether the forgone treatment is low or high value.

\subsection*{Interpretation and scope}
Moral hazard denotes a behavioral response to coverage, not a normative assertion that every additional service is wasteful. With deductibles and caps, a current service can change future marginal prices; estimation must specify patients' expectations and timing. Price sensitivity correlated with risk creates selection on moral hazard, so a common utilization elasticity need not characterize contract choice or optimal design. Treatment-specific cost sharing is constrained by verifiable diagnoses, coding incentives, and administrative expense. Random insurance offers identify intention-to-treat effects; recovering coverage or treatment effects additionally requires compliance assumptions and a well-defined instrument. The extension joins heterogeneity, clinical value, selection, and implementability, recognizing that insurance uncertainty and heterogeneous utilization already have substantial existing results.

\subsection*{Source audit and status}
All three original sources are verified. [S1] supplies foundational uncertainty/information analysis, [S2] randomized cost-sharing evidence, and [S3] selection-on-moral-hazard evidence. The original question therefore cannot be represented as an unanswered claim that such heterogeneity exists. No cited source certifies the proposed joint menu-design problem's current unresolved status.

\subsection*{References}
\begin{enumerate}[label={[S\arabic*]},leftmargin=*,itemsep=0.6em]
\item Kenneth J. Arrow (1963). Uncertainty and the Welfare Economics of Medical Care. American Economic Review 53(5), 941--973.
\url{https://assets.aeaweb.org/asset-server/files/9442.pdf}
\newline\emph{Locator and support limits:} Introduction and article: uncertainty/information foundations; not a treatment-specific insurance optimum.
\item Willard G. Manning, Joseph P. Newhouse, Naihua Duan, Emmett B. Keeler, Arleen Leibowitz, and M. Susan Marquis (1987). Health Insurance and the Demand for Medical Care: Evidence from a Randomized Experiment. American Economic Review 77(3), 251--277.
\url{https://www.jstor.org/stable/1804094}
\newline\emph{Locator and support limits:} Article and RAND experiment: cost-sharing utilization effects; catalogue truncates the author list, while the manuscript includes Marquis.
\item Liran Einav, Amy Finkelstein, Stephen P. Ryan, Paul Schrimpf, and Mark R. Cullen (2013). Selection on Moral Hazard in Health Insurance. American Economic Review 103(1), 178--219. DOI: 10.1257/aer.103.1.178.
\url{https://www.aeaweb.org/articles?id=10.1257/aer.103.1.178}
\newline\emph{Locator and support limits:} Abstract and model: selection on utilization response; health value by treatment is an additional target.
\end{enumerate}

\subsection*{Common conventions: Source mathematical conventions}

Unless an entry states otherwise, agent and firm sets are finite. State and action spaces are measurable, strategies and outcome maps are measurable, probability laws are specified, and expectations exist for the targets under discussion. Infinite-horizon discount factors lie in $(0,1)$ and discounted objectives are finite. Norms, tolerances and complexity input sizes must be specified for computational comparisons. Constants or primitives that appear locally are inputs, not empirical facts asserted by this catalogue.

\textbf{Identification.} A statistical model $m\in\mathcal M$ implies an observable law $P_m$ and target $\theta(m)$. At law $P$, its identified set is
\[
\Theta(P;\mathcal M)=\{\theta(m):m\in\mathcal M,\ P_m=P\}.
\]
Point identification holds when this set is a singleton, or equivalently when $P_{m_1}=P_{m_2}$ implies $\theta(m_1)=\theta(m_2)$ for all compatible models. Sharp bounds mean bounds attained, or approached when the boundary is not attained, by compatible models. Writing this set is a definition, not a solution: the substantive task is to establish informative restrictions and derive or estimate the set. An unrestricted-confounding model may give only trivial bounds.

\textbf{Inference.} Identification, estimability and valid uncertainty quantification are separate questions. Fix $\alpha\in(0,1)$. A confidence procedure $C_n$ over a declared law class $\mathcal P$ has asymptotic uniform coverage at level $1-\alpha$ when
\[
\liminf_{n\to\infty}\inf_{P\in\mathcal P}\Pr_P\{\theta(P)\in C_n\}\geq1-\alpha.
\]
For set-identified targets the coverage event must be defined separately, for example coverage of the full identified set. If an entry uses identification-indexed models instead of uniquely indexed laws, the same coverage convention is applied over those models. Pointwise limits do not establish uniform validity, and asymptotic validity does not guarantee finite-sample precision. Sample sizes, dependence structures and weak-identification sequences are part of each application.

\textbf{Counterfactuals.} $Y_i(a)$ is a potential outcome under a specified intervention, including its timing, implementation and versions. It is not $Y_i$ conditional on voluntarily choosing $a$. A full intervention vector permits interference; shorthand scalar treatment notation does not silently impose no interference. Assignment, selection, attrition, anticipation and measurement must be specified. A claim about an instrument's local effect is not automatically a population policy effect.

\textbf{Economic equilibrium.} A counterfactual model includes best responses, constraints, feasibility, market clearing, state transitions and expectations consistent with the stated information. When several equilibria exist, either a declared selector is an input or the target is set-valued. Comparisons across policy regimes must use the stated selection convention. Reduced-form relationships are not presumed invariant after an equilibrium-changing intervention.

\textbf{Optimization and welfare.} Optimization is over the stated feasible set. If attainment is not established, $\sup$ or $\inf$ and $\varepsilon$-optimal choices replace an asserted maximizer or minimizer. For example, with value $V=\sup_{a\in\mathcal A}W(a)<\infty$, an $\varepsilon$-optimal set is $\{a\in\mathcal A:W(a)\geq V-\varepsilon\}$ for $\varepsilon>0$. Benefits, costs, risk and health or environmental outcomes can be aggregated only using declared units and conversion weights. Interpersonal utility weights, the treatment of fiscal transfers and foreign profits, discounting, and the behavioral welfare criterion are explicit inputs. A comparative-static derivative additionally requires existence and appropriate differentiability of the selected counterfactual.

\textbf{Sources.} Local labels [S1], [S2], and so forth restart in every question. Locators identify the relevant abstract, model, section or result at the level actually checked; a bibliographic or abstract-level verification is not represented as inspection of an entire proof. Working papers are distinguished from later publications and changing versions. Source summaries are paraphrased. All URL checks and editorial formulations refer to the audit date stated above.
% END ECONOMICS STATEMENT
\end{document}
